\documentclass[12pt,a4paper,oneside]{article}
\usepackage[brazil]{babel}
\usepackage[utf8]{inputenc}
\usepackage{olymp}
\usepackage{color}
\usepackage[dvipsnames]{xcolor}
\usepackage{listings}

\setcounter{problem}{3}

\begin{document}

\begin{problem}{Intercalando}{intercala}{2}
 
Faça um programa que receba duas listas de $N$ valores cada e mostre a intercalação dos $2N$ valores das listas.

%\Notes
%
%Observações, se tiver.

\InputFile

A entrada é formada por três linhas: a primeira contém um valor inteiro $N$, indicando a quantidade de números em cada lista; a segunda contém os $N$ números da primeira lista, que são números inteiros $A_i$ separados por espaço em branco; e a terceira linha contém os $N$ números da segunda lista, que são números inteiros $B_i$ separados por espaço em branco.

Restrições: $1 \leq N \leq 1000$, $-10^6 \leq A_i, B_i \leq 10^6$.



\OutputFile

Seu programa deve gerar apenas uma linha de saída contendo $2N$ valores inteiros, representando os elementos intercalados das listas, ou seja, $A_1 \,B_1\, A_2\, B_2\, \dots\, A_n\, B_n$.

Imprima um espaço em branco após cada elemento, inclusive após o último, e não se esqueça de encerrar a linha.

%\Notes
%Nesta questão, seu programa deve usar a função \texttt{fatorial} dada %acima exatamente como está, não a modifique! Sua
%tarefa é apenas criar o programa com a função \texttt{main}.

\Example

\begin{example}
\exmp{
5
0 2 4 6 8
1 3 5 7 9
}{
0 1 2 3 4 5 6 7 8 9
}%
\end{example}

\begin{example}
\exmp{
7
1 1 1 1 1 1 1
2 2 2 2 2 2 2
}{
1 2 1 2 1 2 1 2 1 2 1 2 1 2
}%
\end{example}

%Comentários sobre o exemplo, se necessário.

\end{problem}

\end{document}
